\chapter[Sermon 89. What Shall Be Done When the Thousand Years Are Expired, of the World Deceived, of War and Grievous Persecution of the Godly, and of the Everlasting Pain of the Wicked.]{What Shall Be Done When the Thousand Years Are Expired, of the World Deceived, of War and Grievous Persecution of the Godly, and of the Everlasting Pain of the Wicked.}
\label{sermon:089}
\markright{Sermon 89. What Shall Be Done When the Thousand Years Are ...}

And when the thousand years are expired, Satan shall be loosed out of his prison, and shall go out to deceive the nations that are in the four quarters of the earth, Gog and Magog to gather them together to battle, whose number is as the sand of the sea. And they went upon the plain of the earth, and compassed the tents of the saints about, and the beloved city. And fire came down from God out of heaven and devoured them. And the Devil that deceived them was cast into a lake of fire and brimstone, where the beast and the false prophet are and shall be tormented day and night forevermore.

He declares here what will happen after those thousand years. And he states primarily two things: that the devil will be released from his prison, so that he may deceive the people in the world, and may assemble Gog and Magog for battle. To this he adds two more: a most severe persecution of the church, and punishment of the wicked, and everlasting damnation of the devil and his followers.

\index{John, Saint}And the seduction of the world must again be explained by the figure of synecdoche. For the meaning of the scripture will not allow us to understand that there should be no godly people at that time. For we believe that there is a church, and an holy church, and that it shall always exist in the world until the judgment. And we have heard moreover in this book how many thousands are sealed so that they should not perish. And also that the dragon must be loosed for a little season. Just as we read in the gospel that Satan is cast out and his kingdom taken from him, nevertheless Saint Peter warns and says that the devil goes about like a roaring lion, and seeks whom he may devour. Verily, for the greatest force of Satan is to infringe the faithful, by Christ that mighty champion and noble conqueror, the Devil notwithstanding going about and aspiring again to the empire, and to be restored to his former place. So at this present we understand that Satan, loosed after those thousand years, ranges now abroad more freely, exercises greater authority, seduces more people in the world, and rules further than he has reigned these thousand years. Yet so that there shall nevertheless in the world be a fellowship of Saints dispersed and vexed miserably. For immediately Saint John says that the beloved city of God is besieged by the enemies. Therefore the church shall be in the midst of the enemies. Wherefore all that same place must be explained not of the truth and religion wholly extinguished, but of the more large and ample power and seduction of Satan, the old serpent.

\index{Antichrist}\index{Ezekiel (Book of)}Therefore he says that when the thousand years shall be expired, the Devil will be released from that prison into which, through the power and might of Christ or the preaching of the Apostles, he had been shut. For when the bonds are broken—namely, the sincere doctrine and preaching of the gospel corrupted and depraved—he comes out. And for this purpose he comes out, that he might deceive the Gentiles, that is to say, all people and nations dwelling in the four quarters of the earth, meaning the whole universal world. And to the end he might allure Gog and Magog, namely fierce men, barbarians, worldly, mocking and despising true religion, addicted to robbery and given to evil things, and regarding only corruption and wickedness, that I say, he might draw such men to unrighteousness and keep them still in error. For such does Ezekiel signify Gog and Magog to be. But those who, through divine grace, are not such, shall not be deceived by Satan; but grounded in Christ, shall persevere in the doctrine of the prophets and Apostles, and shall rightly worship Christ, shall abhor Antichrist, and all wickedness in the world.

\index{Jerome, Saint}\index{Mahomet}\index{Turks}\index{Papacy}\index{Daniel (Book of)}\index{Rome}But a deceptive falsehood has passed through the world far and wide, since the thousand years expired, as experience teaches, and histories of times testify. For it is plain that during those thousand years, there were famous churches of Christ in the East, which nevertheless we lament to have been destroyed within these five hundred years. Therefore, the wicked and abominable sect of Mahomet began six hundred years after the birth of Christ, and from that time forth was advanced by the Saracens, but prevailed at last after those thousand fatal years. For how great is the power of the Turks now in Africa, Asia, and Europe, no man is ignorant. And Papistry had its beginning and progression rather soon; but after a thousand years it was of full force. For the Bishops of Rome through the abuse of excommunication have oppressed even the most mighty Emperors and Kings. For who knows not with what shameless boldness the popes have withstood Kings and Emperors, Henrys, Fredericks, Louis, and many other Princes, whom their lewdness has vexed, vanquished, and overcome? After much and grievous contention, the Popes extorted to themselves the consecrating of Bishops. They usurped moreover the church goods also, by which (such a force has money) they might do in the world what they listed. For by this means Papistry received strongest sinews. Moreover, after those thousand years was raised up and established that God Maxim, of whom also Daniel makes mention, which brought also a great strength unto Papistry. I mean transubstantiation, and the horrible polluting of the Lord’s Supper, and manifold abuse of the holy mysteries. And of the force hereof increased an infinite number of priests and filthy Friars. For after those thousand years at the length came up the sect or order of Jacobines, Celestines, Gilbertines, of Gray friars, black friars, white friars, and many other friars, and monstrous Monks, which have craftily crept into the favor of all princes, to the intent they might know all their secrets by auricular confession. Then began all things more impudently to be set forth and sold in the church than ever before. Superstitions and unprofitable and hurtful ceremonies overflowed. For we have seen thirty years since and more, how much increased daily idols and idolatry, worshiping of creatures, and abuses innumerable about the same, pilgrimages to dumb idols, and an infinite number of the same sort. I recite not that holy matrimony waxed now vile after those thousand years, in so much that ministers of churches were prohibited to marry. Then waxed whoredom rife, rape, and adultery, and yet more filthy things than all these, \&c. I pass over here very many things; this only I rehearse, if you compare the rites, ceremonies, and superstitions of Papistry with the heathen gentility (as I have partly showed here and there in my works), you will say that Papistry passes far all gentility. For in case the false opinion and persuasion once taken away, you do way what Papistry is in itself: you will grant that there was never such a corrupt thing in the world. Full rightly therefore says Jerome, that Satan is broken loose out of prison. By which proverb he signifies matters extremely corrupted, nothing to be done in his place or decent order, but all things confused, all turned up side down, at the will and lust of the evil spirit.

Hereunto is added another thing, that when the thousand years expire, Satan should gather Gog and Magog to battle. By these words, John has undoubtedly alluded the prophecy of Ezekiel, which we read in chapters 38 and 39. Ezekiel seems to have prophesied of the wars of Macedon and of Antiochus, speaking of them with a prophetic phrase and a hyperbolic amplification. The prophet says that Gog is the land of Magog. And it is evident that Magog was Japheth’s son, who dwelt at Mount Caucacus and extended his empire to Ethiopia and Egypt. And afterward, out of Asia and from the eastern parts, Antiochus Epiphanes made war on the people of God. This was a figure of Antichrist, as all commentators confess. Therefore, it appears that John brings forth these things by way of comparison, as though he should say: just as in past times the people of Gog and Magog severely molested and afflicted the people of God, so in the times of Antichrist, most grievous wars shall arise, with which the church of God shall be shaken and laid waste. He says truly that the host of these destroyers shall be innumerable. He adds, in the manner of Scripture, a parable for clarity: as the sand of the sea. And also, by another phrase of speaking, he signifies that the enemies of God’s people shall be bold and ready to overrun the whole world and to turmoil all things with wars. For he says: And they went upon the plain of the land. This means that, being swift and bold, they shall run over the whole world; everywhere and throughout the wide world, there shall be cruel wars.

For most purposely he adds: and they surrounded the tents of the Saints and the beloved City. And means that the church of God shall be most grievously plagued with those Gogicall and barbarous wars. For in times past Jerusalem was called the chosen and beloved City; but after she rejected the word of the Lord, she was no more beloved of God, but rather rejected and hated. Therefore Saint John speaks of the Catholic church, which Saint Paul also in another place out of Isaiah names, Jerusalem that is above. The same is also called the tents of Saints. For the faithful are in the church as it were in tents, fighting against Satan, the world, sin, and flesh. And where he says, they surrounded the tents of Saints, he says something more than if he had written, they assailed or besieged or assaulted the tents of Saints. For they surround them, which give the assault round about, and vex them most grievously, as though they were already taken, that no hope can appear to any man, no refuge or way to escape.

Without a doubt, if we compare these things with histories, we will find that the church has been many times assailed with cruel wars; but never yet with crueler than after those thousand fatal years. I mean the holy war, as they term it. William Archbishop of Tyre has written at large about it, as has the Abbot of Uršpurge in Chronicles. Also Benedictus Cottes and Paulus Aemilius in the fourth book on the events of the Franks. Finally, Volaterranus in the eleventh book of Geography, concerning Coele-Syria and Palestine.

Historians report many things of the battle of Troy. Others suppose that those of Assyria and Babylon were greater. Many extol the wars of the Persians and Macedonians, as in truth they were terrible. The Romans also have their Punic, Mithridatic, Civil, Cimbric, and Germanic wars: but I believe that the war they call Holy was more cruel than all these, more bloody and grievous, and of longer duration. In this war, great battles were joined, with countless numbers of men, involving almost all nations and peoples of the world. Wonderful and monstrous slaughters were made. More than a hundred thousand men died, a number beyond belief. It continued for many years, more than any previous war in the world. Furthermore, it was waged with the most hostile intentions. And what is most relevant to this discussion is that the Oriental Saracens, Turks, Egyptians, Babylonians, and other barbarian nations were inflamed in this war, harboring an unquenchable hatred against the Christian religion, seeking to uproot it, and succeeding in uprooting a large portion of it, and they continue to do so to this day. This war, therefore, was the most grievous of all, and was the cause of the persecution of the faithful in the East and West. To note something of this, and to recount for those ignorant of history, it is clear that under Pope Gregory VII, there were many and famous churches in the East, and Patriarchal churches still secure. But while this Pope, above all others, dealt wickedly against Christ, the Son of God, and his holy church, just as we read in the time of Solomon, who, after he had revolted, brought forth many enemies against him, and most cruel ones: so in the wicked and tyrannical reign of Gregory VII, Solyman the Turk invaded Antioch, at which time the Emperors of Greece were said to have been dispatched from the East. And the Turks, advancing, are said to have invaded and harassed first the straits or ports of the Caspian hills and the country of Armenia, around the year of our Lord 764. There is now no time to speak of this in detail. After Solyman, Belchiaroke, the Turkish prince, whom others call Belzet, succeeded him, and he also invaded Greece itself, despising the Emperors of Constantinople. Alexius, who then was Emperor, is said to have sought aid from the Westerners against the Turks. And also one Peter, an Hermit (who certain historians blame most grievously, and not without cause), coming out of the East and traveling through the West, cried out an alarm. Urban II, whom some call Turbane, and a disciple of Gregory VII, called a great council at Clermont in France, where he proposed the question of recovering the Holy Land and delivering the Lord’s sepulcher from the hands of the Infidels. That council reminds me of what is described in the eighth book of Kings, chapter 22, under Ahab and Jehoshaphat, concerning the recovery of Ramoth-Gilead from the hands of the Syrians. For there was also a deceiving spirit there, there were Ahabs, there were Jehoshaphats, and many other similar things. To avoid making many words, a journey was decreed against the barbarian infidels of the East. This was done in the year of our Lord 1095. In the meantime, Peter the Hermit stirred himself and gathered certain thousands, whom he led through Hungary into Asia. Immediately after, followed the unfortunate captains Folkemar and Gottschalk, priests, who, by the way, destroyed everything with fire and sword, and were slain. At last, Godfrey and Baldwin, most noble princes, with certain excellent captains and noble warriors, with an innumerable multitude of men, were transported into Asia: which they say was done in the year of our Lord 1096. And within four years, or at most three, they had taken by assault or surrender the cities of Nice, Heraclea, Tarsus, Antioch, and Jerusalem. The Abbot of Urspurg reports that so much blood was shed in the city of Jerusalem that in the very temple itself, the horses stood up to the knees in the blood of the slain. The same man tells of a notable battle fought at Ascalon, in which about fifteen thousand foot soldiers and five thousand horsemen of Christians overthrew and routed Solyman of Babylon, equipped with one hundred thousand horsemen and four hundred thousand foot soldiers, and that more than one hundred thousand men were slain in that battle. And this journey of Godfrey was the first among the worthy voyages of Syria or Asia.

Following this expedition came others, better equipped. For while the victory and good fortune of those who first went to the East were greatly praised and commended throughout the West, William, Prince and Duke of Poitiers, filled with great hope, led approximately one hundred thousand soldiers into the Eastern lands. The year of our Lord was reckoned 1101. But of so great a number, scarcely one thousand are recorded to have returned home safely.

3. After, in the year of our Lord 1147, through the urging of Bernard of Clairvaux, Louis King of France, and Conrad King of Germany, and Frederick Prince of Swabia, took their journey into the East, leading with them an army almost innumerable; but the same army perished almost entirely, with scarcely the princes surviving.

In the year of our Lord 1189, when the city of Jerusalem was taken by the Sultan, the King of Persia—a city the Christians had held for only about 89 years—Emperor Frederick, nicknamed Barbarossa, King Philip of France, King Richard of England, and other powerful princes raised a very large army of Christian people to recover the city and the Holy Land. They were fortunate to transport their army into Asia, but afterward suffered great misfortune. For Emperor Frederick was drowned, and the entire army, as Urſpurgenſ testifies, died of the plague.

5. The fifteenth (and that famous indeed) voyage in Syria involved the mighty Kings Philip of France and Richard of England, nicknamed “Lionheart.” This occurred in the year of our Lord 1191. However, they returned without accomplishing anything of significance, and with a considerable loss of men.

6. And Palmerius, a chronicler, says that Henry reports the son of Emperor Barbarossa sent an army into Syria, which returned the following year. Consequently, the Christians, lacking assistance in Syria, lost all the dominion they had retained. He speaks of these things in the year of our Lord 1198.

7 Again, in the year of our Lord 1213, Pope Innocent III of that name sent public letters to all the faithful of Christ, urging them to take up arms against the infidels who possessed the Holy Land. If any man has leisure and desires to read the letters, he will find them in the Chronicles of Ursperg. And then, in the year of our Lord 1215, he held a general council in Lateran, where war was decreed against the Easterlings. Also, Honorius III, about the year of our Lord 1217, addressed and confirmed the same thing. Consequently, many Christian princes met at Acre, formerly called Ptolemais, and waged a deadly war upon the Easterlings. In this conflict, they captured the noble city of Damietta. Yet neither the outcome nor the fruit justified such great enterprises, costs, perils, and losses.

8 Therefore Frederick II, the emperor, hoping to do some good, marches also with an eager and well-equipped army into the East; this they say occurred in the year of our Lord 1234. In the meantime, while he valiantly fights in the East, the Bishop of Rome, Gregory IX of that name, taking an occasion (I use the words of Ursperger) of the emperor’s absence, sent a great army into Apulia and seized the emperor’s lands. This was done while the emperor was absent serving Christ (which is most wicked to say), and kept them subdued for his own use, and by no means would allow those who had taken the holy cross (that is to say, those who were going to war for the emperor) to take shipping or passage, but kept them under his power, both in Apulia and in Lombardy. And more such matters, which those with leisure may read in the same [? chronicle]. Therefore, the emperor, constrained, leaving his affairs there unfinished, was forced to negotiate with the enemy and returned, that he might recover what the Pope had taken from him.

And not long afterward, namely in the year of our Lord 1248, Louis, King of France, with his brothers Robert and Charles, and a very powerful army, sailed into Syria. There Robert was slain, and Charles was captured by the Sultan, and was scarcely released at the last, with great difficulty.

King Louis of France, in the year of our Lord 1270, embarked with his three sons from Marseille to sail into Africa. A plague struck his army in enemy territory, by which both the father and the sons died, and the entire army suffered an exceedingly great calamity.

And again, although they experienced misfortune in the wars against the Barbarians, it was nevertheless considered again in the council of Lyons under Gregory X, about the year of our Lord 1273, concerning the recovery of the holy land. But Palmerius in the year of our Lord 1291 states that many thousands of Christians were slain in Syria by the Saracens, and all the rest fled the country in fear. And the Chronicle of the Kings of France states that Aemilius ended the holy war (namely, in the year of our Lord 1291) when Ptolemais in the East was destroyed by the Sultan. It is manifest therefore that this barbaric and gogical war has lasted about 195 years—such a long time as I know no other war in the world that was ever waged with such obstinate minds, with so great armies, and so much shedding of human blood. We see in the meantime the tents of saints, and the city of God beloved, namely, the faithful church throughout the world, especially in the East, and also in the West, to be most grievously afflicted, and more than oppressed and destroyed, with only a few small remnants remaining; so that not without cause we may perceive what the Lord said in the gospel: but when the Son of Man shall come, will he find any faith on Earth?

\index{Judgment, Last}The most holy and wise Prophet of God Daniel seems to have foreseen and prophesied all those things, as he did all the rest concerning Antichrist. After speaking at length about the power of Antichrist and the worship of God against the Apostles' institution, he adds in the 11th Chapter. And in the time of the end, to wit the end of the world and last judgment approaching, shall set upon him, namely upon Antichrist, the king of the South, and the king of the North shall fall upon him like a whirlwind, with chariots and horsemen, with a strong and great navy, and shall invade his realms. He shall overflow with armies, to wit innumerable, and he shall pass through, that is to say, he shall overcome all like a conqueror, doing what he lists. For we have perceived that the armies sent into the East by the counsels and motion of the Bishop of Rome have molested by sea and land the Turks and also the Soldan of Babylon and Egypt. What will they say that Daniel, pointing as it were with his finger, adds? He shall come also into the chosen land and invade the land of desire, namely Judea, which some time was called the chosen, delectable, and pleasant land. And many shall fall in the war verily, that shall be made for the recovery of the holy land. It follows in Daniel, these shall be delivered out of his hand, Aedom and Moab, and the Princes of the children of Ammon. For those nations are not known to have been so destroyed as the rest were, by the Saracens and after by the Turks, for that they framed themselves to them in time. Daniel annexes, and he shall lay his hands upon realms, neither shall the land of Egypt escape. For it is evident that the same also was possessed of the Soldans, princes of Babylon, and of the emperors of the Turks. It follows, and he shall have the rule of the treasures of gold and silver and all the precious things of the Egyptians. By which the prophet has signified the inestimable treasures and riches, and excellent majesty of the Soldans and Turkish Emperors. All which things, even so as the Prophet has said, experience proves to have been and as yet to be fulfilled. The Prophet adds, finally, the Libyans and Aethiopians shall be in his journeys. This translator has turned, “He shall pass also through Lybia and Aethiopia,” or as others have translated it, “they shall be in his way.” And he means that those regions shall be open to those barbarous Soldans and emperors of the Turks, by league, vicinity, and amity. Jerome expounding this place says, “when Egypt was taken, those lands were also afraid.” Wherefore he says not that he took them, but passed through Lybia and Aethiopia. Whether sense of these soever thou choosest, thou shalt not err, as I think, from the truth. And Daniel adds the disturbance from the East and from the North shall trouble him, so much that he shall go forth in a great fury to destroy and kill many. Jerome shows must be understood of Antichrist. The Pope of Rome affirms that the Patriarchal seats are subject to him, as Jerusalem, Antioch and Alexandria, and the holy land to be his right. And he hears from the East and from the North that all those parts are possessed of the Soldans and Emperors of Turks; he calls therefore great councils and decrees war against them. He hears moreover that Constantinople is taken, that the Rhodes is won, Dalmatia subdued, Bulgaria and Hungary vanquished, and so on. Again therefore he summoneth councils, he arms kings, he leads forth soldiers, he moves war, and decrees that war shall be made for the recovering of the holy land and to root out the Turks. So verily this Gogmagog war is not yet ended or appeased at this day. Whereby it comes to pass that an infinite multitude of men are slain on either side. Furthermore at the end of this Prophecy, the prophet shows, and as it were with his finger, the palace or seat of Antichrist, by Antiochus figured before, lest any man should not know where Antichrist were to be found. And he shall plant, or set the tabernacle of his palace betwixt two Seas: to wit the Adriatic Sea, called now the gulf of Venice, and the Tyrrhenian or Tuscan Sea, in the mount of desire of holiness, that is to say in the pleasant and holy hill. We have heard certainly that the palace of Saint Peter is preferred both before Mount Zion and also Sinai. There sits the most holy, in the seat of holiness. There is most full remission of all sins. There is the mother and supreme head of all churches. There is the high court and judgment, from whence may no man appeal. There sits the king of kings and high Bishop, which so far excels in brightness and majesty the Emperor and other kings, as the sun doth the moon and stars. There is thought to be perfect holiness, and all the treasures of Christ and of his Saints. Therefore said Daniel rightly, that Antichrist shall dwell in the noble and holy hill, namely in the seven hilly Rome, as we heard also in the 17th Chapter. Finally, he prophesies also of the end of this most puissant prince, Antichrist, and says: and what time he shall come to his end, no man shall help him. For Christ coming to judgment, shall thrust him out of his seat. And Daniel in the 12th chapter following describes the judgment. To Christ alone be glory.

Let us therefore proceed to add a few things concerning the torments of the ungodly, and the everlasting condemnation of the Devil and his followers. John says, “And fire came down from heaven and devoured them.” And the prophet Amos, in chapter one, calls God’s vengeance fire, as others do also. Therefore, John signifies that the vengeance of God will fall upon all the enemies of the church. In times past, fire coming down from heaven burned up Sodom and Gomorrah, and also consumed the enemies of Elijah. And although fire does not always fall from heaven in a physical sense, the persecutors of the church will never escape punishment, because they have afflicted Christ’s saints. Surely, if we observe what happened in that holy war, and what happens daily, we will say that God’s vengeance is most present both against the Turks and the Papists. But if anyone understands that about the end of the world fire will rage and consume the wicked, as Peter also mentions of fire and burning, 2 Peter 3, I will not contradict it.

Finally, he also addresses the everlasting damnation of Satan and all his followers. For where the Lord said in the Gospel, “If the blind lead the blind, both shall fall into the ditch,” it follows that both Satan, the deceiver, and the people he has seduced, will be carried together to hell. There, John places, as it were, joining together the devil, Gog and Magog, the Saracens, Turks, briefly all deceived nations, the Beast, the false prophet, and all Antichristians. We see therefore that God’s judgment is righteous, which he now returns to describe. And we admonished before, by this speech, that they shall be tormented day and night, and so forth, to signify the perpetuity of damnation. From this, may the Lord our God deliver us: to whom be glory forevermore. Amen.
